\documentclass[11pt]{article}
\newcommand{\LabNumber}{5}
\newcommand{\LabTitle}{Keypad Scanning and Display}
\newcommand{\LabTopic}{Multiplexed keypad input scanning on the Arduino Uno R3, decoded and shown on a single 7-segment digit in AVR assembly -- the input-side counterpart to Lab 4's multiplexed display output}
\newcommand{\LabDuration}{3 hours (in-lab)}
% Shared preamble for Computer Programming lab sheets.
% Include with % Shared preamble for Computer Programming lab sheets.
% Include with % Shared preamble for Computer Programming lab sheets.
% Include with \input{common.tex} after \documentclass{article}.

\usepackage[a4paper,margin=2.2cm]{geometry}
\usepackage[T1]{fontenc}
\usepackage{lmodern}
\usepackage{microtype}
\usepackage{enumitem}
\usepackage{xcolor}
\usepackage{tcolorbox}
\tcbuselibrary{breakable,skins}
\usepackage{listings}
\usepackage{fancyhdr}
\usepackage{titlesec}
\usepackage{amsmath}
\usepackage{array}
\usepackage{booktabs}
\usepackage{hyperref}

\hypersetup{
  colorlinks=true,
  linkcolor=black,
  urlcolor=blue!60!black,
  pdfborder={0 0 0}
}

% ---------- Course identity (override per file if needed) ----------
\providecommand{\CourseName}{Microprocessors}
\providecommand{\CourseCode}{010153521}
\providecommand{\Semester}{Semester 1/2026}
\providecommand{\LabDuration}{3 hours (in-lab)}

% ---------- Colors ----------
\definecolor{accent}{HTML}{1F4E79}
\definecolor{accentlight}{HTML}{D9E2EC}
\definecolor{warn}{HTML}{B45309}
\definecolor{ok}{HTML}{166534}
\definecolor{codebg}{HTML}{F6F8FA}
\definecolor{codeframe}{HTML}{D0D7DE}
\definecolor{codekw}{HTML}{0550AE}
\definecolor{codestr}{HTML}{0A7B2C}
\definecolor{codecom}{HTML}{6E7781}

% ---------- Code / assembly listings ----------
\lstdefinestyle{c}{
  language=C,
  basicstyle=\ttfamily\small,
  keywordstyle=\color{codekw}\bfseries,
  stringstyle=\color{codestr},
  commentstyle=\color{codecom}\itshape,
  backgroundcolor=\color{codebg},
  frame=single,
  rulecolor=\color{codeframe},
  framerule=0.4pt,
  showstringspaces=false,
  breaklines=true,
  columns=fullflexible,
  keepspaces=true,
  tabsize=4,
  xleftmargin=8pt,
  xrightmargin=4pt,
  aboveskip=6pt,
  belowskip=6pt,
}
\lstdefinestyle{asm}{
  basicstyle=\ttfamily\small,
  commentstyle=\color{codecom}\itshape,
  backgroundcolor=\color{codebg},
  frame=single,
  rulecolor=\color{codeframe},
  framerule=0.4pt,
  showstringspaces=false,
  breaklines=true,
  columns=fullflexible,
  keepspaces=true,
  tabsize=4,
  xleftmargin=8pt,
  xrightmargin=4pt,
  aboveskip=6pt,
  belowskip=6pt,
}
\lstset{style=asm}

% ---------- Section styling ----------
\titleformat{\section}{\Large\bfseries\color{accent}}{\thesection.}{0.5em}{}
\titleformat{\subsection}{\large\bfseries\color{accent!85!black}}{\thesubsection}{0.5em}{}
\titlespacing*{\section}{0pt}{12pt}{6pt}
\titlespacing*{\subsection}{0pt}{8pt}{4pt}

% ---------- Page headers/footers ----------
\pagestyle{fancy}
\fancyhf{}
\fancyhead[L]{\small\CourseName\ifx\CourseCode\empty\else\ (\CourseCode)\fi}
\fancyhead[C]{\small\Semester}
\fancyhead[R]{\small\textbf{Lab \LabNumber:} \LabTitle}
\fancyfoot[C]{\small\thepage}
\renewcommand{\headrulewidth}{0.4pt}

% ---------- Title block ----------
\newcommand{\labheader}{%
  \begin{tcolorbox}[
    enhanced,
    colback=accentlight,
    colframe=accent,
    boxrule=0.5pt,
    arc=2pt,
    left=10pt, right=10pt, top=8pt, bottom=8pt
  ]
    {\Large\bfseries\color{accent} Lab \LabNumber: \LabTitle}\\[2pt]
    {\small \CourseName\ifx\CourseCode\empty\else\ (\CourseCode)\fi\quad\textbar\quad \Semester}\\[1pt]
    {\small \textbf{Duration:} \LabDuration\quad\textbar\quad \textbf{Topic:} \LabTopic}
  \end{tcolorbox}
  \vspace{2pt}
}

% ---------- Custom environments ----------
\newtcolorbox{summary}[1][Quick Reference]{
  enhanced,breakable,
  colback=accentlight!40,colframe=accent!70,
  title={\bfseries #1},
  fonttitle=\color{white},
  attach boxed title to top left={xshift=6pt,yshift=-3pt},
  boxed title style={colback=accent,boxrule=0pt,arc=2pt},
  arc=2pt,boxrule=0.4pt,
  top=12pt,
}

\newcounter{labex}[section]
\newtcolorbox{labexercise}[1]{
  enhanced,breakable,
  colback=white,colframe=accent,
  title={\bfseries In-Lab Exercise \refstepcounter{labex}\thelabex: #1},
  fonttitle=\color{white},
  attach boxed title to top left={xshift=6pt,yshift=-3pt},
  boxed title style={colback=accent,boxrule=0pt,arc=2pt},
  arc=2pt,boxrule=0.4pt,
  top=12pt,
}

\newcounter{traceex}
\newtcolorbox{trace}[1][]{
  enhanced,breakable,
  colback=white,colframe=warn,
  title={\bfseries Trace \& Predict \refstepcounter{traceex}\thetraceex\ifx\\#1\\\else: #1\fi},
  fonttitle=\color{white},
  attach boxed title to top left={xshift=6pt,yshift=-3pt},
  boxed title style={colback=warn,boxrule=0pt,arc=2pt},
  arc=2pt,boxrule=0.4pt,
  top=12pt,
}

\newcounter{bugex}
\newtcolorbox{bug}[1][]{
  enhanced,breakable,
  colback=white,colframe=warn,
  title={\bfseries Find the Bug \refstepcounter{bugex}\thebugex\ifx\\#1\\\else: #1\fi},
  fonttitle=\color{white},
  attach boxed title to top left={xshift=6pt,yshift=-3pt},
  boxed title style={colback=warn,boxrule=0pt,arc=2pt},
  arc=2pt,boxrule=0.4pt,
  top=12pt,
}

\newcounter{writeex}
\newtcolorbox{writecode}[1][]{
  enhanced,breakable,
  colback=white,colframe=warn,
  title={\bfseries Write Code on Paper \refstepcounter{writeex}\thewriteex\ifx\\#1\\\else: #1\fi},
  fonttitle=\color{white},
  attach boxed title to top left={xshift=6pt,yshift=-3pt},
  boxed title style={colback=warn,boxrule=0pt,arc=2pt},
  arc=2pt,boxrule=0.4pt,
  top=12pt,
}

% ---------- AI Collaboration box ----------
\definecolor{aipurple}{HTML}{4C1D95}
\definecolor{aipurplelight}{HTML}{EDE9FE}

\newtcolorbox{aicollab}[1][AI Collaboration]{
  enhanced,breakable,
  colback=aipurplelight,colframe=aipurple!70,
  title={\bfseries #1},
  fonttitle=\color{white},
  attach boxed title to top left={xshift=6pt,yshift=-3pt},
  boxed title style={colback=aipurple,boxrule=0pt,arc=2pt},
  arc=2pt,boxrule=0.4pt,
  top=12pt,
}

\newcounter{tryex}
\newtcolorbox{tryit}[1]{
  enhanced,breakable,
  colback=ok!6,colframe=ok!70!black,
  title={\bfseries Try It \refstepcounter{tryex}\thetryex: #1},
  fonttitle=\color{white},
  attach boxed title to top left={xshift=6pt,yshift=-3pt},
  boxed title style={colback=ok!70!black,boxrule=0pt,arc=2pt},
  arc=2pt,boxrule=0.4pt,
  top=12pt,
}

% ---------- Drawing exercise ----------
\definecolor{drawteal}{HTML}{0D7377}
\definecolor{drawteallight}{HTML}{D4F1F4}

\newcounter{drawex}
\newtcolorbox{drawex}[1]{
  enhanced,breakable,
  colback=drawteallight,colframe=drawteal!70,
  title={\bfseries Draw on Paper \refstepcounter{drawex}\thedrawex: #1},
  fonttitle=\color{white},
  attach boxed title to top left={xshift=6pt,yshift=-3pt},
  boxed title style={colback=drawteal,boxrule=0pt,arc=2pt},
  arc=2pt,boxrule=0.4pt,
  top=12pt,
}


\newcommand{\takehomebanner}{%
  \vspace{6pt}
  \begin{tcolorbox}[
    colback=warn!10,colframe=warn,boxrule=0.5pt,arc=2pt,
    left=10pt,right=10pt,top=6pt,bottom=6pt
  ]
    \textbf{Take-Home (Paper Work).}\;
    Solve the following on paper without using a computer or compiler.
    Hand-write your answers; submit at the start of the next lab.
    Goal: prepare you to dry-code during the written exam.
  \end{tcolorbox}
  \vspace{4pt}
}
 after \documentclass{article}.

\usepackage[a4paper,margin=2.2cm]{geometry}
\usepackage[T1]{fontenc}
\usepackage{lmodern}
\usepackage{microtype}
\usepackage{enumitem}
\usepackage{xcolor}
\usepackage{tcolorbox}
\tcbuselibrary{breakable,skins}
\usepackage{listings}
\usepackage{fancyhdr}
\usepackage{titlesec}
\usepackage{amsmath}
\usepackage{array}
\usepackage{booktabs}
\usepackage{hyperref}

\hypersetup{
  colorlinks=true,
  linkcolor=black,
  urlcolor=blue!60!black,
  pdfborder={0 0 0}
}

% ---------- Course identity (override per file if needed) ----------
\providecommand{\CourseName}{Microprocessors}
\providecommand{\CourseCode}{010153521}
\providecommand{\Semester}{Semester 1/2026}
\providecommand{\LabDuration}{3 hours (in-lab)}

% ---------- Colors ----------
\definecolor{accent}{HTML}{1F4E79}
\definecolor{accentlight}{HTML}{D9E2EC}
\definecolor{warn}{HTML}{B45309}
\definecolor{ok}{HTML}{166534}
\definecolor{codebg}{HTML}{F6F8FA}
\definecolor{codeframe}{HTML}{D0D7DE}
\definecolor{codekw}{HTML}{0550AE}
\definecolor{codestr}{HTML}{0A7B2C}
\definecolor{codecom}{HTML}{6E7781}

% ---------- Code / assembly listings ----------
\lstdefinestyle{c}{
  language=C,
  basicstyle=\ttfamily\small,
  keywordstyle=\color{codekw}\bfseries,
  stringstyle=\color{codestr},
  commentstyle=\color{codecom}\itshape,
  backgroundcolor=\color{codebg},
  frame=single,
  rulecolor=\color{codeframe},
  framerule=0.4pt,
  showstringspaces=false,
  breaklines=true,
  columns=fullflexible,
  keepspaces=true,
  tabsize=4,
  xleftmargin=8pt,
  xrightmargin=4pt,
  aboveskip=6pt,
  belowskip=6pt,
}
\lstdefinestyle{asm}{
  basicstyle=\ttfamily\small,
  commentstyle=\color{codecom}\itshape,
  backgroundcolor=\color{codebg},
  frame=single,
  rulecolor=\color{codeframe},
  framerule=0.4pt,
  showstringspaces=false,
  breaklines=true,
  columns=fullflexible,
  keepspaces=true,
  tabsize=4,
  xleftmargin=8pt,
  xrightmargin=4pt,
  aboveskip=6pt,
  belowskip=6pt,
}
\lstset{style=asm}

% ---------- Section styling ----------
\titleformat{\section}{\Large\bfseries\color{accent}}{\thesection.}{0.5em}{}
\titleformat{\subsection}{\large\bfseries\color{accent!85!black}}{\thesubsection}{0.5em}{}
\titlespacing*{\section}{0pt}{12pt}{6pt}
\titlespacing*{\subsection}{0pt}{8pt}{4pt}

% ---------- Page headers/footers ----------
\pagestyle{fancy}
\fancyhf{}
\fancyhead[L]{\small\CourseName\ifx\CourseCode\empty\else\ (\CourseCode)\fi}
\fancyhead[C]{\small\Semester}
\fancyhead[R]{\small\textbf{Lab \LabNumber:} \LabTitle}
\fancyfoot[C]{\small\thepage}
\renewcommand{\headrulewidth}{0.4pt}

% ---------- Title block ----------
\newcommand{\labheader}{%
  \begin{tcolorbox}[
    enhanced,
    colback=accentlight,
    colframe=accent,
    boxrule=0.5pt,
    arc=2pt,
    left=10pt, right=10pt, top=8pt, bottom=8pt
  ]
    {\Large\bfseries\color{accent} Lab \LabNumber: \LabTitle}\\[2pt]
    {\small \CourseName\ifx\CourseCode\empty\else\ (\CourseCode)\fi\quad\textbar\quad \Semester}\\[1pt]
    {\small \textbf{Duration:} \LabDuration\quad\textbar\quad \textbf{Topic:} \LabTopic}
  \end{tcolorbox}
  \vspace{2pt}
}

% ---------- Custom environments ----------
\newtcolorbox{summary}[1][Quick Reference]{
  enhanced,breakable,
  colback=accentlight!40,colframe=accent!70,
  title={\bfseries #1},
  fonttitle=\color{white},
  attach boxed title to top left={xshift=6pt,yshift=-3pt},
  boxed title style={colback=accent,boxrule=0pt,arc=2pt},
  arc=2pt,boxrule=0.4pt,
  top=12pt,
}

\newcounter{labex}[section]
\newtcolorbox{labexercise}[1]{
  enhanced,breakable,
  colback=white,colframe=accent,
  title={\bfseries In-Lab Exercise \refstepcounter{labex}\thelabex: #1},
  fonttitle=\color{white},
  attach boxed title to top left={xshift=6pt,yshift=-3pt},
  boxed title style={colback=accent,boxrule=0pt,arc=2pt},
  arc=2pt,boxrule=0.4pt,
  top=12pt,
}

\newcounter{traceex}
\newtcolorbox{trace}[1][]{
  enhanced,breakable,
  colback=white,colframe=warn,
  title={\bfseries Trace \& Predict \refstepcounter{traceex}\thetraceex\ifx\\#1\\\else: #1\fi},
  fonttitle=\color{white},
  attach boxed title to top left={xshift=6pt,yshift=-3pt},
  boxed title style={colback=warn,boxrule=0pt,arc=2pt},
  arc=2pt,boxrule=0.4pt,
  top=12pt,
}

\newcounter{bugex}
\newtcolorbox{bug}[1][]{
  enhanced,breakable,
  colback=white,colframe=warn,
  title={\bfseries Find the Bug \refstepcounter{bugex}\thebugex\ifx\\#1\\\else: #1\fi},
  fonttitle=\color{white},
  attach boxed title to top left={xshift=6pt,yshift=-3pt},
  boxed title style={colback=warn,boxrule=0pt,arc=2pt},
  arc=2pt,boxrule=0.4pt,
  top=12pt,
}

\newcounter{writeex}
\newtcolorbox{writecode}[1][]{
  enhanced,breakable,
  colback=white,colframe=warn,
  title={\bfseries Write Code on Paper \refstepcounter{writeex}\thewriteex\ifx\\#1\\\else: #1\fi},
  fonttitle=\color{white},
  attach boxed title to top left={xshift=6pt,yshift=-3pt},
  boxed title style={colback=warn,boxrule=0pt,arc=2pt},
  arc=2pt,boxrule=0.4pt,
  top=12pt,
}

% ---------- AI Collaboration box ----------
\definecolor{aipurple}{HTML}{4C1D95}
\definecolor{aipurplelight}{HTML}{EDE9FE}

\newtcolorbox{aicollab}[1][AI Collaboration]{
  enhanced,breakable,
  colback=aipurplelight,colframe=aipurple!70,
  title={\bfseries #1},
  fonttitle=\color{white},
  attach boxed title to top left={xshift=6pt,yshift=-3pt},
  boxed title style={colback=aipurple,boxrule=0pt,arc=2pt},
  arc=2pt,boxrule=0.4pt,
  top=12pt,
}

\newcounter{tryex}
\newtcolorbox{tryit}[1]{
  enhanced,breakable,
  colback=ok!6,colframe=ok!70!black,
  title={\bfseries Try It \refstepcounter{tryex}\thetryex: #1},
  fonttitle=\color{white},
  attach boxed title to top left={xshift=6pt,yshift=-3pt},
  boxed title style={colback=ok!70!black,boxrule=0pt,arc=2pt},
  arc=2pt,boxrule=0.4pt,
  top=12pt,
}

% ---------- Drawing exercise ----------
\definecolor{drawteal}{HTML}{0D7377}
\definecolor{drawteallight}{HTML}{D4F1F4}

\newcounter{drawex}
\newtcolorbox{drawex}[1]{
  enhanced,breakable,
  colback=drawteallight,colframe=drawteal!70,
  title={\bfseries Draw on Paper \refstepcounter{drawex}\thedrawex: #1},
  fonttitle=\color{white},
  attach boxed title to top left={xshift=6pt,yshift=-3pt},
  boxed title style={colback=drawteal,boxrule=0pt,arc=2pt},
  arc=2pt,boxrule=0.4pt,
  top=12pt,
}


\newcommand{\takehomebanner}{%
  \vspace{6pt}
  \begin{tcolorbox}[
    colback=warn!10,colframe=warn,boxrule=0.5pt,arc=2pt,
    left=10pt,right=10pt,top=6pt,bottom=6pt
  ]
    \textbf{Take-Home (Paper Work).}\;
    Solve the following on paper without using a computer or compiler.
    Hand-write your answers; submit at the start of the next lab.
    Goal: prepare you to dry-code during the written exam.
  \end{tcolorbox}
  \vspace{4pt}
}
 after \documentclass{article}.

\usepackage[a4paper,margin=2.2cm]{geometry}
\usepackage[T1]{fontenc}
\usepackage{lmodern}
\usepackage{microtype}
\usepackage{enumitem}
\usepackage{xcolor}
\usepackage{tcolorbox}
\tcbuselibrary{breakable,skins}
\usepackage{listings}
\usepackage{fancyhdr}
\usepackage{titlesec}
\usepackage{amsmath}
\usepackage{array}
\usepackage{booktabs}
\usepackage{hyperref}

\hypersetup{
  colorlinks=true,
  linkcolor=black,
  urlcolor=blue!60!black,
  pdfborder={0 0 0}
}

% ---------- Course identity (override per file if needed) ----------
\providecommand{\CourseName}{Microprocessors}
\providecommand{\CourseCode}{010153521}
\providecommand{\Semester}{Semester 1/2026}
\providecommand{\LabDuration}{3 hours (in-lab)}

% ---------- Colors ----------
\definecolor{accent}{HTML}{1F4E79}
\definecolor{accentlight}{HTML}{D9E2EC}
\definecolor{warn}{HTML}{B45309}
\definecolor{ok}{HTML}{166534}
\definecolor{codebg}{HTML}{F6F8FA}
\definecolor{codeframe}{HTML}{D0D7DE}
\definecolor{codekw}{HTML}{0550AE}
\definecolor{codestr}{HTML}{0A7B2C}
\definecolor{codecom}{HTML}{6E7781}

% ---------- Code / assembly listings ----------
\lstdefinestyle{c}{
  language=C,
  basicstyle=\ttfamily\small,
  keywordstyle=\color{codekw}\bfseries,
  stringstyle=\color{codestr},
  commentstyle=\color{codecom}\itshape,
  backgroundcolor=\color{codebg},
  frame=single,
  rulecolor=\color{codeframe},
  framerule=0.4pt,
  showstringspaces=false,
  breaklines=true,
  columns=fullflexible,
  keepspaces=true,
  tabsize=4,
  xleftmargin=8pt,
  xrightmargin=4pt,
  aboveskip=6pt,
  belowskip=6pt,
}
\lstdefinestyle{asm}{
  basicstyle=\ttfamily\small,
  commentstyle=\color{codecom}\itshape,
  backgroundcolor=\color{codebg},
  frame=single,
  rulecolor=\color{codeframe},
  framerule=0.4pt,
  showstringspaces=false,
  breaklines=true,
  columns=fullflexible,
  keepspaces=true,
  tabsize=4,
  xleftmargin=8pt,
  xrightmargin=4pt,
  aboveskip=6pt,
  belowskip=6pt,
}
\lstset{style=asm}

% ---------- Section styling ----------
\titleformat{\section}{\Large\bfseries\color{accent}}{\thesection.}{0.5em}{}
\titleformat{\subsection}{\large\bfseries\color{accent!85!black}}{\thesubsection}{0.5em}{}
\titlespacing*{\section}{0pt}{12pt}{6pt}
\titlespacing*{\subsection}{0pt}{8pt}{4pt}

% ---------- Page headers/footers ----------
\pagestyle{fancy}
\fancyhf{}
\fancyhead[L]{\small\CourseName\ifx\CourseCode\empty\else\ (\CourseCode)\fi}
\fancyhead[C]{\small\Semester}
\fancyhead[R]{\small\textbf{Lab \LabNumber:} \LabTitle}
\fancyfoot[C]{\small\thepage}
\renewcommand{\headrulewidth}{0.4pt}

% ---------- Title block ----------
\newcommand{\labheader}{%
  \begin{tcolorbox}[
    enhanced,
    colback=accentlight,
    colframe=accent,
    boxrule=0.5pt,
    arc=2pt,
    left=10pt, right=10pt, top=8pt, bottom=8pt
  ]
    {\Large\bfseries\color{accent} Lab \LabNumber: \LabTitle}\\[2pt]
    {\small \CourseName\ifx\CourseCode\empty\else\ (\CourseCode)\fi\quad\textbar\quad \Semester}\\[1pt]
    {\small \textbf{Duration:} \LabDuration\quad\textbar\quad \textbf{Topic:} \LabTopic}
  \end{tcolorbox}
  \vspace{2pt}
}

% ---------- Custom environments ----------
\newtcolorbox{summary}[1][Quick Reference]{
  enhanced,breakable,
  colback=accentlight!40,colframe=accent!70,
  title={\bfseries #1},
  fonttitle=\color{white},
  attach boxed title to top left={xshift=6pt,yshift=-3pt},
  boxed title style={colback=accent,boxrule=0pt,arc=2pt},
  arc=2pt,boxrule=0.4pt,
  top=12pt,
}

\newcounter{labex}[section]
\newtcolorbox{labexercise}[1]{
  enhanced,breakable,
  colback=white,colframe=accent,
  title={\bfseries In-Lab Exercise \refstepcounter{labex}\thelabex: #1},
  fonttitle=\color{white},
  attach boxed title to top left={xshift=6pt,yshift=-3pt},
  boxed title style={colback=accent,boxrule=0pt,arc=2pt},
  arc=2pt,boxrule=0.4pt,
  top=12pt,
}

\newcounter{traceex}
\newtcolorbox{trace}[1][]{
  enhanced,breakable,
  colback=white,colframe=warn,
  title={\bfseries Trace \& Predict \refstepcounter{traceex}\thetraceex\ifx\\#1\\\else: #1\fi},
  fonttitle=\color{white},
  attach boxed title to top left={xshift=6pt,yshift=-3pt},
  boxed title style={colback=warn,boxrule=0pt,arc=2pt},
  arc=2pt,boxrule=0.4pt,
  top=12pt,
}

\newcounter{bugex}
\newtcolorbox{bug}[1][]{
  enhanced,breakable,
  colback=white,colframe=warn,
  title={\bfseries Find the Bug \refstepcounter{bugex}\thebugex\ifx\\#1\\\else: #1\fi},
  fonttitle=\color{white},
  attach boxed title to top left={xshift=6pt,yshift=-3pt},
  boxed title style={colback=warn,boxrule=0pt,arc=2pt},
  arc=2pt,boxrule=0.4pt,
  top=12pt,
}

\newcounter{writeex}
\newtcolorbox{writecode}[1][]{
  enhanced,breakable,
  colback=white,colframe=warn,
  title={\bfseries Write Code on Paper \refstepcounter{writeex}\thewriteex\ifx\\#1\\\else: #1\fi},
  fonttitle=\color{white},
  attach boxed title to top left={xshift=6pt,yshift=-3pt},
  boxed title style={colback=warn,boxrule=0pt,arc=2pt},
  arc=2pt,boxrule=0.4pt,
  top=12pt,
}

% ---------- AI Collaboration box ----------
\definecolor{aipurple}{HTML}{4C1D95}
\definecolor{aipurplelight}{HTML}{EDE9FE}

\newtcolorbox{aicollab}[1][AI Collaboration]{
  enhanced,breakable,
  colback=aipurplelight,colframe=aipurple!70,
  title={\bfseries #1},
  fonttitle=\color{white},
  attach boxed title to top left={xshift=6pt,yshift=-3pt},
  boxed title style={colback=aipurple,boxrule=0pt,arc=2pt},
  arc=2pt,boxrule=0.4pt,
  top=12pt,
}

\newcounter{tryex}
\newtcolorbox{tryit}[1]{
  enhanced,breakable,
  colback=ok!6,colframe=ok!70!black,
  title={\bfseries Try It \refstepcounter{tryex}\thetryex: #1},
  fonttitle=\color{white},
  attach boxed title to top left={xshift=6pt,yshift=-3pt},
  boxed title style={colback=ok!70!black,boxrule=0pt,arc=2pt},
  arc=2pt,boxrule=0.4pt,
  top=12pt,
}

% ---------- Drawing exercise ----------
\definecolor{drawteal}{HTML}{0D7377}
\definecolor{drawteallight}{HTML}{D4F1F4}

\newcounter{drawex}
\newtcolorbox{drawex}[1]{
  enhanced,breakable,
  colback=drawteallight,colframe=drawteal!70,
  title={\bfseries Draw on Paper \refstepcounter{drawex}\thedrawex: #1},
  fonttitle=\color{white},
  attach boxed title to top left={xshift=6pt,yshift=-3pt},
  boxed title style={colback=drawteal,boxrule=0pt,arc=2pt},
  arc=2pt,boxrule=0.4pt,
  top=12pt,
}


\newcommand{\takehomebanner}{%
  \vspace{6pt}
  \begin{tcolorbox}[
    colback=warn!10,colframe=warn,boxrule=0.5pt,arc=2pt,
    left=10pt,right=10pt,top=6pt,bottom=6pt
  ]
    \textbf{Take-Home (Paper Work).}\;
    Solve the following on paper without using a computer or compiler.
    Hand-write your answers; submit at the start of the next lab.
    Goal: prepare you to dry-code during the written exam.
  \end{tcolorbox}
  \vspace{4pt}
}

\usepackage{circuitikz}

\begin{document}
\labheader

\section*{Objectives}
\begin{itemize}[leftmargin=*,nosep]
  \item Understand \textbf{matrix keypad scanning} as multiplexed \emph{input} --
        the mirror image of Lab 4's multiplexed \emph{output} to a 2-digit display.
  \item Wire a 4x4 matrix keypad and a single 7-segment digit to the Arduino Uno R3.
  \item Write AVR assembly that drives one row \textbf{LOW} at a time, polls all
        four columns (internal pull-ups, idle \textbf{HIGH}), and decodes the
        resulting (row, column) pair into a key value.
  \item Debounce keypresses and display the result -- reusing Lab 4-1's
        \texttt{DIGIT\_TO\_SEG} subroutine \emph{unchanged}.
\end{itemize}

\begin{summary}[Quick Reference: Pin Assignment]
\begin{center}
\begin{tabular}{@{}p{2.1cm}p{1.8cm}p{2.6cm}p{7.4cm}@{}}
\toprule
\textbf{Port bits} & \textbf{Uno pins} & \textbf{Direction} & \textbf{Function} \\
\midrule
\texttt{PB0-PB3} & D8-D11 & Output & Keypad rows \texttt{R1-R4} (driven one at a
time, active \textbf{low}) \\
\texttt{PD2-PD5} & D2-D5 & Input, pull-up & Keypad columns \texttt{C1-C4} (idle
\textbf{high}; pulled \textbf{low} through a pressed key when its row is driven low) \\
\texttt{PC0-PC5} & A0-A5 & Output & Segments \texttt{a,b,c,d,e,f} \\
\texttt{PD6} & D6 & Output & Segment \texttt{g} \\
\bottomrule
\end{tabular}
\end{center}
This is the \emph{same} segment bit order and hex table as Lab 4-1 -- only segment
\texttt{g}'s physical pin moved, from \texttt{PD2} to \texttt{PD6}, because
\texttt{PD2-PD5} are now the keypad's column inputs. There is no digit-select
transistor this time: a single digit's common cathode goes straight to \textbf{GND}.
\end{summary}

\begin{summary}[Quick Reference: How Matrix Scanning Works]
A 4x4 keypad packs 16 keys into just 8 wires -- 4 rows + 4 columns -- because each
key does nothing but bridge \emph{one} row wire to \emph{one} column wire when
pressed; there is no other connection inside the keypad. Driving all 4 rows at once
would make every key on a pressed column look identical, so scanning removes the
ambiguity: drive exactly \emph{one} row \textbf{LOW} (the other three stay
\textbf{HIGH}), read all 4 columns, and only trust what you read while that one row
is the low one. Cycle through all 4 rows, fast enough that a keypress during any
row's turn gets caught.
\begin{center}
\begin{tabular}{@{}c|cccc@{}}
\toprule
 & \textbf{C1} & \textbf{C2} & \textbf{C3} & \textbf{C4} \\
\midrule
\textbf{R1} & 1 & 2 & 3 & A \\
\textbf{R2} & 4 & 5 & 6 & B \\
\textbf{R3} & 7 & 8 & 9 & C \\
\textbf{R4} & * & 0 & \# & D \\
\bottomrule
\end{tabular}
\end{center}
This is the standard layout for most inexpensive 4x4 membrane keypads --
\textbf{confirm it against your physical keypad's printed labels} before trusting
any decoded output; wiring order varies between manufacturers. Only \texttt{0-9}
have a segment code from Lab 4-1's table; for \texttt{A,B,C,D,*,\#}, just display a
dash (segment \texttt{g} only) as a ``not a digit'' indicator -- no new segment
codes needed.
\end{summary}

\begin{summary}[Quick Reference: Your Individual Target]
As in Labs 3 and 4-1, derive your personal target from your student ID:
\[
D = \text{sum of every digit in your student ID}, \qquad
t_{\text{debounce}} = 15\text{ ms} + (D \bmod 10) \times 1\text{ ms}
\]
\textbf{Worked example} (a made-up ID, do not reuse): ID \texttt{1029384756123}
$\to D = 51$. Then $t_{\text{debounce}} = 15 + (51 \bmod 10) \times 1 = 15 + 1 =
16$~ms.
\end{summary}

\section*{Bill of Materials}
\begin{center}
\begin{tabular}{@{}p{1.4cm}p{6.3cm}p{7.5cm}@{}}
\toprule
\textbf{Qty} & \textbf{Part} & \textbf{Notes} \\
\midrule
1 & Arduino Uno R3 (from Lab 3/4) & Supplies power, ground, and all signal pins \\
1 & 4x4 matrix keypad module (8-pin) & \texttt{R1-R4}, \texttt{C1-C4} \\
1 & 7-segment display, \textbf{common cathode} & Single digit -- no transistor
needed \\
7 & Resistor, 330~$\Omega$ & Segment current limiting, \texttt{a-g} \\
1 & Breadboard + jumper wires & \\
\bottomrule
\end{tabular}
\end{center}

\section*{Building the Circuit}
\begin{enumerate}[leftmargin=*,nosep]
  \item \textbf{Rows.} Wire keypad \texttt{R1-R4} to \texttt{PB0-PB3} (Uno D8-D11).
        These are \textbf{outputs} -- the MCU actively drives them, one at a time.
  \item \textbf{Columns.} Wire keypad \texttt{C1-C4} to \texttt{PD2-PD5} (Uno
        D2-D5). These are \textbf{inputs} with the AVR's internal pull-up enabled
        -- no external resistor, exactly like Lab 4-1's buttons.
  \item \textbf{Display.} Wire segments \texttt{a-f} through 330~$\Omega$ resistors
        to \texttt{PC0-PC5} (A0-A5), segment \texttt{g} through a 330~$\Omega$
        resistor to \texttt{PD6} (D6), and the common cathode straight to
        \textbf{GND}. One digit, always lit when driven -- no multiplexed refresh
        loop this time.
\end{enumerate}

\begin{figure}[h!]
\centering
\begin{circuitikz}[
  american voltages,
  scale=0.85, transform shape,
  font=\scriptsize
]

% ---------------- Keypad box ----------------
\node[draw, fill=accentlight, minimum width=3cm, minimum height=3cm, align=center]
     at (3,5) (kp) {\textbf{4x4 KEYPAD}};

% Row pins (left edge) <- PB0-PB3
\draw (0,6.0) node[left]{\texttt{PB0} (D8)}  -- (1.5,6.0) node[pos=0.85,above]{R1};
\draw (0,5.5) node[left]{\texttt{PB1} (D9)}  -- (1.5,5.5) node[pos=0.85,above]{R2};
\draw (0,5.0) node[left]{\texttt{PB2} (D10)} -- (1.5,5.0) node[pos=0.85,above]{R3};
\draw (0,4.5) node[left]{\texttt{PB3} (D11)} -- (1.5,4.5) node[pos=0.85,above]{R4};

% Column pins (right edge) -> PD2-PD5
\draw (4.5,6.0) -- (6,6.0) node[right]{\texttt{PD2} (D2)};
\node[above] at (4.725,6.0) {C1};
\draw (4.5,5.5) -- (6,5.5) node[right]{\texttt{PD3} (D3)};
\node[above] at (4.725,5.5) {C2};
\draw (4.5,5.0) -- (6,5.0) node[right]{\texttt{PD4} (D4)};
\node[above] at (4.725,5.0) {C3};
\draw (4.5,4.5) -- (6,4.5) node[right]{\texttt{PD5} (D5)};
\node[above] at (4.725,4.5) {C4};

% ---------------- Segment section ----------------
\node[draw, fill=accentlight, minimum width=3cm, minimum height=1.6cm, align=center]
     at (3,0.5) (dig) {\textbf{DIGIT}\\common cathode};

\node[align=center] at (3,3.4) {\scriptsize \texttt{PC0-PC5} (A0-A5), \texttt{PD6} (D6)};

\foreach \i/\lbl in {0/a, 1/b, 2/c, 3/d, 4/e, 5/f, 6/g} {
  \draw (1.65+\i*0.47,3.0) to[R] (1.65+\i*0.47,1.3);
  \node[below right] at (1.65+\i*0.47,1.35) {\tiny \lbl};
}

\draw (dig.south) -- ++(0,-0.8) node[ground]{};

\end{circuitikz}
\caption{Keypad and single-digit display wiring. Each key merely bridges its row
line to its column line when pressed -- there is no component inside the keypad
beyond that mechanical bridge. Segment resistors shown individually since, unlike
Lab 4-1, there is only one digit to drive (no shared bus, no multiplexing).}
\end{figure}

\section*{Quick Experiments (Try It)}

\begin{tryit}{Rows and columns, by hand}
Before writing any scanning code, drive \texttt{PB0} (Row 1) \textbf{LOW} and
\texttt{PB1-PB3} \textbf{HIGH}, with columns configured as instructed above. Touch
a jumper wire between the \texttt{R1} and \texttt{C1} pins on the breadboard
(simulating pressing the ``1'' key) and confirm \texttt{IN R17,PIND} shows bit~2
(\texttt{PD2}, column 1) go \textbf{LOW} only while the wire bridges those two
specific pins. This confirms your wiring before any scanning loop exists.
\end{tryit}

\begin{tryit}{One row, no full scan}
Write a short loop that drives \emph{only} \texttt{R1} low (\texttt{R2-R4} stay
high) and polls \texttt{C1-C4}. Confirm that keys \texttt{1, 2, 3, A} register
while keys on other rows (e.g.\ \texttt{5} or \texttt{9}) do \emph{not} -- even
though you have not yet written code to scan those other rows. This is your
checkpoint that a single row's scan works in isolation, the same incremental
approach Lab 4-1 used before combining two digits.
\end{tryit}

\section*{In-Lab Exercises (target: 3 hours)}

\begin{labexercise}{Wire and verify the display alone \hfill {\small \itshape (30 min)}}
Wire the display per the instructions above and re-run Lab 4-1's static-digit test
(cycle 0-9, holding each briefly) \emph{completely unchanged} except for segment
\texttt{g}'s pin (\texttt{PD2} $\to$ \texttt{PD6}). Confirm every digit 0-9 still
renders correctly before wiring the keypad -- you are only relocating one pin, not
redesigning the display logic.
\end{labexercise}

\begin{labexercise}{Single-row scan \hfill {\small \itshape (35 min)}}
Repeat Try It 2 as a graded checkpoint: scan \texttt{R1} only, and light the
Uno's on-board LED (\texttt{PB5}) -- or show it on the display as a placeholder --
whenever any of \texttt{C1-C4} reads low. Confirm the LED responds \emph{only} to
keys \texttt{1, 2, 3, A} and stays off for every other key on the pad.
\end{labexercise}

\begin{labexercise}{Full 4x4 scan and decode \hfill {\small \itshape (90 min)}}
Extend Exercise 2 into a loop that cycles all four rows. On detecting a pressed
key, decode its (row, column) pair into a key value using the Quick Reference
map above, and display it: \texttt{0-9} through \texttt{DIGIT\_TO\_SEG} (reused
unchanged from Lab 4-1), anything else as a dash (segment \texttt{g} only).
\begin{itemize}[nosep,leftmargin=*]
  \item Add a software debounce using your $t_{\text{debounce}}$ from the
        ``Your Individual Target'' box: after first detecting a key, wait
        $t_{\text{debounce}}$ and confirm the same key still reads pressed before
        accepting it.
  \item \textbf{Trace/predict first, then verify:} before running it, predict what
        your program displays if two keys in \emph{different} rows and
        \emph{different} columns are held down at the same time (e.g.\ \texttt{2}
        and \texttt{8}). Then try it and compare -- does your prediction match,
        and can you explain any mismatch from how the scan visits rows in order?
\end{itemize}
\end{labexercise}

\begin{labexercise}{Hold last key until the next press \hfill {\small \itshape (25 min)}}
Confirm your debounce value matches your calculated $t_{\text{debounce}}$ exactly,
then extend Exercise 3 so the display holds the \emph{last} key pressed
persistently -- not just for an instant -- until a new key is pressed and
debounced. Describe in your notebook what the display does differently with
versus without the debounce check from Exercise 3 (e.g.\ does the digit ever
flicker or jump to a wrong value on a real keypress?).
\end{labexercise}

\begin{aicollab}
\textbf{Try these prompts with an AI tool:}
\begin{itemize}[leftmargin=*,nosep]
  \item ``How do I scan a 4x4 matrix keypad in AVR assembly using internal
        pull-up resistors on the columns?''
  \item ``Why does a cheap matrix keypad without diodes suffer from `ghost'
        keypresses when multiple keys are held down at once?''
  \item ``What's a clean way in assembly to convert a (row, column) pair into a
        single key value, given the keys aren't in a simple row-times-four-plus-
        column order on every keypad?''
\end{itemize}

\medskip\textbf{Critical check -- verify, don't just accept:}
Before trusting any decoded output, verify your \emph{physical} keypad's actual
row/column/key layout against the Quick Reference table above -- don't assume an
AI tool's row/column map matches your specific keypad, since wiring order varies
between manufacturers even for keypads that look identical.

\medskip\textbf{Know the limit:}
AI tools will often hand you scanning code that quietly assumes a diode-protected
keypad capable of full N-key rollover. Most classroom kits are diode-less, so
holding multiple keys at once can produce a ``phantom'' key that was never
pressed -- a real hardware limitation the AI's example code won't warn you about
unless you ask specifically.
\end{aicollab}

\takehomebanner

\begin{trace}[Decode a scan cycle by hand]
A scan cycle drives \texttt{R1, R2, R3, R4} low in turn. While \texttt{R3} is low,
\texttt{IN R17,PIND} reads \texttt{0b11110111} (bits 2-5 are columns \texttt{C1-C4}
in that order, bit 2 = C1). Every other row's read shows all four column bits
high. Using the Quick Reference key map, which key was pressed?
\end{trace}

\begin{bug}[Pull-ups left disabled]
A student wired the columns to \texttt{PD2-PD5} exactly as instructed, correctly
set \texttt{DDRD} bits 2-5 to \texttt{0} (input), but forgot to set the
corresponding \texttt{PORTD} bits to \texttt{1} -- so the internal pull-ups never
turn on. No key is being pressed. Predict what \texttt{IN R17,PIND} reads on
these four bits, and explain why an unconnected (floating) digital input does not
reliably read as either a clean \texttt{0} or a clean \texttt{1}.
\end{bug}

\begin{writecode}[Hand-write the decode chain]
On paper, without a computer, simulator, or AI tool, write the \texttt{CPI}/\texttt{BRNE}
comparison chain (or equivalent) that decodes a scan result for \emph{row R2 only}
into its four key values (\texttt{4, 5, 6, B}) -- i.e.\ given that R2 was the row
just driven low, and one column bit is known to read low, write the branches that
turn ``column 1 is low'' into displaying \texttt{4}, ``column 2 is low'' into
\texttt{5}, and so on.
\end{writecode}

\begin{drawex}{Draw the internal matrix}
Draw the keypad's internal 4x4 grid from memory: 4 horizontal row lines, 4
vertical column lines, and a switch symbol at each of the 16 crossings, labeled
with its key. Then, for the key \texttt{6}, trace and highlight the exact current
path from \texttt{PB1} (R2, driven low) through the pressed switch to \texttt{PD4}
(C3, read low) -- and explain, from your drawing, why no \emph{other} column
reads low while only the \texttt{6} key is held.
\end{drawex}

\section*{Reflection}

In one short paragraph, compare this lab's keypad scanning to Lab 4's display
multiplexing. Both drive one line at a time in a repeating cycle -- but one is
scanning \emph{output} (lighting a digit) and the other is scanning \emph{input}
(reading a key). What stayed structurally the same between them (the one-at-a-time
timing, the loop-over-four-lines shape), and what's fundamentally different (which
direction data flows, and why this lab's scan rate has no flicker-fusion deadline
the way Lab 4's display refresh did)?

\end{document}
